\section{Boolean representation and explicit normal forms}
\label{ca-normal}

This section studies classical propositional formulas as finite syntactic
objects inside an intuitionistic arithmetic.  Every use of a truth value below
means a bit, and every choice is a search through an explicitly bounded finite
list.  In particular, a case distinction on the value of a formula is a case
distinction on a primitive recursive Boolean computation; it does not assert
excluded middle for arbitrary arithmetical propositions.

Fix an ordered list $X=(p_0,\ldots,p_{n-1})$ containing the variables of the
formulas under discussion.  An assignment is a number $a<2^n$, its $i$th bit
being the value of $p_i$.  Write $A[a]\in\{0,1\}$ for Boolean evaluation.
The notation $A\equiv_2 B$ abbreviates
\[
  \bigwedge_{a<2^n}\bigl(A[a]=B[a]\bigr).
\]
This is a bounded Boolean test, independent of adding unused variables to $X$.
Empty conjunctions have value $1$ and are written $\top$; empty disjunctions
have value $0$ and are written $\bot$.  A Boolean function on these variables
means an explicitly supplied table of $2^n$ bits, encoded by a number less than
$2^{2^n}$.  This convention avoids quantification over an unrestricted
function space.

Size bounds below refer to finite syntactic data: $n$ is the number of
distinct input variables and $s$ is the number of syntax-tree nodes, after
densely renaming the variables in their order of first occurrence.  The
renaming and its inverse are finite lists and preserve evaluation.
Arithmetic codes for these objects may be much larger than their string
lengths; no complexity statement treats the numerical code as the input
length.  In particular, a gate label below has value less than $n+s$ and
requires $O(\log(n+s+2))$ binary digits.

\subsection{Cofactors, canonical forms, and negation normal form}

\begin{theorem}[Shannon expansion and restriction]\label{ca-shannon}
Let $A_0=A[p:=\bot]$ and $A_1=A[p:=\top]$.  Then
\[
 A\equiv_2(\neg p\wedge A_0)\vee(p\wedge A_1)
 \equiv_2(p\vee A_0)\wedge(\neg p\vee A_1).
\]
Substitution of a bit for $p$ commutes with evaluation.  Consequently these
expansions remain valid when $p$ is replaced by an arbitrary formula and
the corresponding cofactors are substituted uniformly.
\end{theorem}
\begin{proof}
The substitution assertion follows by recursion on the syntax: a variable
equal to $p$ takes the substituted value; another variable is unchanged; and
each connective applies the same Boolean operation to the recursively
computed values of its arguments.  For an assignment with $p=0$, the first
displayed disjunction becomes $A_0$, and its displayed conjunctive version
also becomes $A_0$.  For $p=1$ both become $A_1$.  These two cases exhaust the
values of the bit representing $p$.  Uniform substitution follows from the
same substitution lemma, with the bit now obtained by evaluating the
substituted formula.  All recursive assertions can be checked simultaneously
on a list of subformulas ordered before their parents.
\end{proof}

For $a<2^n$, let
\[
 \ell_i^a=\begin{cases}p_i&\text{if the $i$th bit of $a$ is $1$},\\
                        \neg p_i&\text{otherwise},\end{cases}
 \qquad M_a=\bigwedge_{i<n}\ell_i^a,
 \qquad K_a=\bigvee_{i<n}\neg\ell_i^a.
\]
Here $\neg\ell_i^a$ can immediately be simplified to the opposite literal.
The minterm $M_a$ is true at exactly $a$; the clause $K_a$ is false at exactly
$a$.  These assertions follow by scanning the $n$ bit positions: equality
of assignments means equality at every position, and inequality supplies a
least differing position by bounded search.

\begin{theorem}[Canonical DNF and CNF]\label{ca-canonical}
For every formula $A$,
\[
 D_A=\bigvee_{\substack{a<2^n\\ A[a]=1}}M_a,
 \qquad
 C_A=\bigwedge_{\substack{a<2^n\\ A[a]=0}}K_a
\]
satisfy $A\equiv_2D_A\equiv_2C_A$.  Each form has at most $2^n$ terms or
clauses and at most $n2^n$ literal occurrences.  The same constructions
represent every supplied Boolean table.  Two formulas have identical
canonical DNF, with a fixed ordering and fixed conventions for empty folds,
exactly when they are Boolean equivalent.
\end{theorem}
\begin{proof}
Fix $b<2^n$.  Among the minterms only $M_b$ is true at $b$.  Thus $D_A[b]=1$
exactly when that minterm was included, namely when $A[b]=1$.  Among the
clauses only $K_b$ is false at $b$.  Hence $C_A[b]=0$ exactly when that clause
was included, namely when $A[b]=0$.  The bounds follow directly from the
outer loop over $2^n$ rows and the inner loop over $n$ positions.  For $n=0$
there is one empty assignment, one minterm $\top$, and one clause $\bot$;
the argument therefore includes constant functions.  Equality of the ordered
lists of selected rows is equivalent to equality of their selection bits.
No appeal to completeness or to an infinite Boolean algebra is involved.
\end{proof}

\begin{theorem}[Effective negation normal form]\label{ca-nnf}
Every formula in $\bot,\top,\neg,\wedge,\vee,\to$ has a Boolean-equivalent
formula in which negation occurs only on variables.  The transformation is
primitive recursive and preserves the set of variables.
\end{theorem}
\begin{proof}
Construct simultaneously $P(A)$ and $N(A)$, representing $A$ and $\neg A$.
For a variable use $p$ and $\neg p$; for $\bot$ use $\bot,\top$ and for
$\top$ use $\top,\bot$.  At a negation swap the two outputs.  At a
conjunction put
\[
 P(B\wedge C)=P(B)\wedge P(C),\qquad
 N(B\wedge C)=N(B)\vee N(C).
\]
At a disjunction interchange $\wedge$ and $\vee$ in these formulas.  Finally
put
\[
 P(B\to C)=N(B)\vee P(C),\qquad
 N(B\to C)=P(B)\wedge N(C).
\]
Each clause follows by examining the two input bits; the assertions for all
nodes follow by induction on the finite syntax list.  Biconditionals, if
present as abbreviations, may first be expanded as a conjunction of two
implications.  The simultaneous construction is a finite bottom-up pass;
duplicating earlier finite outputs gives a primitive recursive tree code.
The negative clauses here are identities of Boolean semantics.  They are
not being asserted as unrestricted equivalences of intuitionistic formulas.
\end{proof}

\subsection{Functional completeness and exact obstructions}

\begin{theorem}[Complete Boolean bases]\label{ca-bases}
The bases $\{\neg,\wedge,\vee,\bot,\top\}$,
$\{\neg,\wedge,\bot,\top\}$, and $\{\neg,\vee,\bot,\top\}$ represent
every finite Boolean table.  With at least one available variable, either
NAND alone or NOR alone also represents every Boolean table, including the
two constant functions.
\end{theorem}
\begin{proof}
The first claim is the DNF construction.  The identities
$x\vee y=\neg(\neg x\wedge\neg y)$ and
$x\wedge y=\neg(\neg x\vee\neg y)$ establish the next two claims by
substitution.  If $x\mathbin{\uparrow}y=\neg(x\wedge y)$, then
\[
 \neg x=x\mathbin{\uparrow}x,\quad
 x\wedge y=(x\mathbin{\uparrow}y)\mathbin{\uparrow}
              (x\mathbin{\uparrow}y),\quad
 x\vee y=(x\mathbin{\uparrow}x)\mathbin{\uparrow}
              (y\mathbin{\uparrow}y).
\]
The term $p\mathbin{\uparrow}(p\mathbin{\uparrow}p)$ is constantly true;
its NAND with itself is constantly false.  For
$x\mathbin{\downarrow}y=\neg(x\vee y)$ the dual identities give negation,
disjunction, and conjunction, and
$p\mathbin{\downarrow}(p\mathbin{\downarrow}p)$ is constantly false.
Recursively replace each connective by the indicated fixed template.
Each replacement duplicates its already constructed arguments a fixed
number of times.  A recurrence $b(0)=20$,
$b(t+1)=20b(t)+20$ bounds the output length for input depth at most $t$,
including the constant templates.  This supplies an explicit primitive
recursive bound even without sharing.  With no variables and no nullary
symbols, there are no closed NAND-only or NOR-only terms; the qualification
about an available variable is therefore necessary.
\end{proof}

\begin{theorem}[Monotone representation]\label{ca-monotone}
A supplied Boolean table $f$ is representable using only
$\wedge,\vee,\bot,\top$ if and only if it is monotone: whenever
$a_i\leq b_i$ for every $i<n$, one has $f(a)\leq f(b)$.
\end{theorem}
\begin{proof}
Constants and coordinate projections are monotone, and conjunction and
disjunction preserve monotonicity.  Syntax induction proves necessity.
For sufficiency define
\[
  F=\bigvee_{\substack{a<2^n\\ f(a)=1}}
             \bigwedge_{\substack{i<n\\ a_i=1}}p_i.
\]
If $f(b)=1$, the term indexed by $b$ makes $F(b)=1$.  If $F(b)=1$, bounded
search finds an included $a$ whose true coordinates all occur among the true
coordinates of $b$.  Thus $a\leq b$ coordinatewise, and monotonicity gives
$f(b)=1$.  This proves equality of the tables.  In particular, the basis
containing only conjunction and disjunction cannot express negation, since
negation reverses the order between the one-bit assignments $0$ and $1$.
\end{proof}

\subsection{Boolean algebra and algebraic normal form}

\begin{theorem}[The finite algebra of formulas]\label{ca-finite-algebra}
Modulo Boolean equivalence on $n$ variables, formulas have exactly
$2^{2^n}$ equivalence classes.  Their Boolean algebra is explicitly
isomorphic to the algebra of all subsets of the $2^n$ assignments, with
intersection, union, complement, the empty set, and the whole set.
Its atoms are the $2^n$ minterms.
\end{theorem}
\begin{proof}
Send a formula to its truth-table bitset.  Conjunction, disjunction, and
negation compute intersection, union, and complement bit by bit; equal
bitsets are exactly equivalent formulas.  Every bitset has a preimage by
canonical DNF.  Thus distinct bitsets give distinct classes, and every class
has exactly one bitset.  A nonempty bitset contains a least set bit, so every
nonzero element contains the singleton at that position.  A singleton has
no nonzero proper subset.  Under the isomorphism these singletons are exactly
the minterms.  Counting bitsets of length $2^n$ gives the stated cardinality.
All statements can be expressed with the canonical bitsets as representatives,
so no quotient type or comprehension over an infinite set is needed in iPRA.
The usual commutative, associative, distributive, absorption, complement,
identity, and De Morgan laws follow from the corresponding one-bit checks.
\end{proof}

There is also a completely finite meaning of the word \emph{free} in this
example.  Supply a finite Boolean algebra by finite operation tables satisfying
the Boolean equations, together with $n$ chosen elements $b_i$.  There is a
unique Boolean homomorphism from the truth-table algebra taking the $i$th
coordinate table to $b_i$.  Indeed, send the singleton indexed by $a$ to
$\bigwedge_{i<n}b_i^{a_i}$, where $b_i^1=b_i$ and $b_i^0$ is its complement,
and send a bitset to the join of its selected singleton images.  The images
are pairwise disjoint: different assignments contribute complementary
factors at a differing position.  Their join is $1$ by distributing
$\bigwedge_i(b_i\vee\neg b_i)$.  These two facts prove preservation of joins,
meets, and complement.  They also show that the coordinate tables have the
required images.  Uniqueness follows because every bitset is the join of its
singletons and each singleton is the indicated expression in the generators.
Every distribution here is a finite induction on $n$ or on a bitset length.

Use $x\oplus y$ for exclusive or and multiplication for conjunction of bits.
Then $x\oplus x=0$, $x^2=x$, and multiplication distributes over exclusive
or.  These are verified by their four input cases.

\begin{theorem}[Unique algebraic normal form]\label{ca-anf}
For every Boolean table $f$ on $n$ variables there are unique coefficients
$c_S\in\{0,1\}$, indexed by subsets $S\subseteq\{0,\ldots,n-1\}$, such that
\[
 f(x)=\bigoplus_{S\subseteq\{0,\ldots,n-1\}}
                  c_S\prod_{i\in S}x_i.
\]
They are given by the finite transform
\[
 c_S=\bigoplus_{T\subseteq S}f(1_T),
\]
where $1_T$ is the assignment with exactly the coordinates in $T$ true.
The coefficient transform is its own inverse.
\end{theorem}
\begin{proof}
At $x=1_U$, only monomials with $S\subseteq U$ contribute.  Substituting
the proposed coefficients gives
\[
 \bigoplus_{S\subseteq U}\ \bigoplus_{T\subseteq S}f(1_T).
\]
For fixed $T\subseteq U$, its contribution occurs once for every $S$
between $T$ and $U$.  If $T=U$ there is one occurrence.  Otherwise choose
the least element of $U\setminus T$; toggling its membership in $S$ pairs
all occurrences without a fixed point.  Equal bits cancel in exclusive or,
leaving exactly $f(1_U)$.  This proves both the representation and that
applying the transform twice returns the original table.  Consequently it
is injective, proving uniqueness.  Alternatively, $c_\varnothing=f(0)$,
and the value at $1_U$ determines $c_U$ after the coefficients for proper
subsets have been determined.  Enumeration of pairs $(T,S)$ is bounded by
$4^n$, so both the construction and the pairing argument are primitive
recursive finite computations.
\end{proof}

\begin{corollary}[Affine fragment]\label{ca-affine}
The functions representable by exclusive or, negation, and constants are
exactly $c\oplus\bigoplus_{i\in S}x_i$.  This basis cannot express
conjunction of two distinct variables.
\end{corollary}
\begin{proof}
Exclusive or combines coefficients modulo $2$, and negation toggles the
constant coefficient, so structural induction proves that every represented
function has this form.  Conversely the displayed form uses precisely those
operations.  The algebraic normal form of $x_0x_1$ has a degree-two monomial;
uniqueness in Theorem~\ref{ca-anf} excludes an affine representation.
\end{proof}

\subsection{Prime implicants and bounded minimization}

A consistent term is a conjunction of literals containing at most one sign
of each variable.  It is coded by a ternary vector of length $n$: absent,
positive, or negative.  It is an \emph{implicant} of $f$ if every assignment
satisfying it makes $f$ true.  It is \emph{prime} if deleting any one of its
literals destroys this property.  Both tests inspect at most $2^n$
assignments for each of at most $n$ deletions.

\begin{theorem}[Complete prime-implicant representation]\label{ca-prime}
The disjunction of all prime implicants of $f$ represents $f$.  There are
at most $3^n$ candidate terms.  Dually, the conjunction of all prime implicates
represents $f$.
\end{theorem}
\begin{proof}
Every included term implies $f$, by definition.  If $f(a)=1$, its complete
minterm is an implicant.  Repeatedly delete the least literal whose deletion
preserves being an implicant.  Each successful deletion decreases the number
of literals, so $n$ iterations, with idle steps after stopping, suffice.
The result remains true at $a$ and permits no further deletion.  If some
proper subterm were an implicant, deleting just one of the removed literals
would also have been an implicant, since a stronger term implies a weaker
one.  Thus the output is prime and covers $a$.  For the dual statement apply
this argument to $1-f$ and negate each term using De Morgan's law.  This also
handles the constantly false function, whose prime-implicant disjunction is
empty, and the constantly true function, whose empty term is prime.
\end{proof}

The theorem supplies a primitive recursive, exact minimization procedure
among prime-implicant covers for any primitive recursive cost function, or
for costs supplied as a finite input table: enumerate the at most
$2^{3^n}$ subsets of the candidate prime implicants, retain the subsets
representing $f$, and choose a least-cost retained subset with a fixed tie
rule.  At least one is retained.  A minimum among arbitrary
consistent-term DNFs may be sought among prime implicants when the cost is
the number of terms, the number of literal occurrences, or a positive
combination of these: replacing each term by a containing prime implicant
does not increase either count and preserves the represented function.
This is a termination and correctness result, not an efficiency claim.

\section{Clause representations and certified satisfiability procedures}
\label{ca-solvers}

A literal is a variable with a sign, its complement being denoted
$\overline\ell$.  A clause is a finite set of literals, interpreted as their
disjunction.  The empty clause is false.  A clause set is interpreted as the
conjunction of its clauses; the empty clause set is true.  Removing repeated
literals or clauses preserves evaluation.  A clause containing a literal
and its complement is always true and may be discarded.  Thus a normalized,
non-tautological clause on $n$ variables has a ternary code, and there are at
most $3^n$ such clauses, including the empty one.

\subsection{A linear-size extension to clauses of width at most three}

Associate a fresh variable $z$ with each non-variable node of a formula's
syntax tree.  Process children before parents, using an original variable
as the label of its own leaf.  For each node impose the following clauses,
where $a,b$ are its child labels:
\[
\begin{array}{c|l}
 \text{node relation}&\text{clauses}\\ \hline
 z=\neg a &(\neg z\vee\neg a),\ (z\vee a)\\
 z=a\wedge b &(\neg z\vee a),\ (\neg z\vee b),\
                   (z\vee\neg a\vee\neg b)\\
 z=a\vee b &(z\vee\neg a),\ (z\vee\neg b),\
                   (\neg z\vee a\vee b)\\
 z=(a\to b)&(\neg z\vee\neg a\vee b),\ (z\vee a),\
                   (z\vee\neg b)\\
 z=\top &(z)\\
 z=\bot &(\neg z).
\end{array}
\]
Add the unit clause containing the root label.  Call the resulting clause
set $T_A$.  Here ``3-CNF'' means that every clause has at most three literal
occurrences.  If exactly three occurrences are required, a shorter clause
can be padded by repeating a literal; no fresh variables are needed for that
convention.

\begin{theorem}[Tseitin extension and projection]\label{ca-tseitin}
For every assignment $a$ to the original variables:
\begin{enumerate}
 \item there is exactly one extension to the fresh labels satisfying all
       node clauses;
 \item that extension satisfies $T_A$ if and only if $A[a]=1$;
 \item every satisfying assignment of $T_A$ restricts to a satisfying
       assignment of $A$.
\end{enumerate}
If $s$ is the number of nodes of $A$, the construction uses at most $s$
fresh variables, at most $3s+1$ clauses, and at most $7s+1$ literal
occurrences.  In particular it preserves the number of satisfying assignments
when assignments are counted over precisely the original variables and the
introduced labels.
\end{theorem}
\begin{proof}
For each row of the table, the conjunction of its clauses is true exactly
when the stated equation of bits holds.  For example, the two clauses
$\neg z\vee a$ and $\neg z\vee b$ say that $z=1$ forces both inputs to be
$1$; the third says that when both inputs are $1$, $z$ must be $1$.  The
other cases are verified by the same finite inspection.
Given $a$, label each node by evaluating its children.  Each node equation
then holds, so the node clauses hold.  Conversely, the bit at a leaf is
already fixed, and each node equation uniquely determines its label from its
children.  Induction along the node order proves that every satisfying
extension has exactly the labels just computed.  Its root label is $A[a]$,
which proves all three claims after adding the root unit clause.  Every
binary gate contributes at most three clauses and seven occurrences; a
negation or constant contributes fewer.  Add one occurrence for the root.
Uniqueness gives a bijection between satisfying original assignments and
satisfying extended assignments, not merely preservation of existence.
\end{proof}

Here nodes are occurrences, so two identical subformulas at different
positions receive distinct fresh labels; their defining equations force the
same bit nevertheless.  Sharing identical subformulas is an optional
optimization, provided the shared directed graph remains acyclic.  With
binary-encoded dense labels, the output bit length is
$O(s\log(n+s+2))$; the displayed linear bounds concern clauses and literal
occurrences.

This is an equivalence after bounded existential elimination of the fresh
labels.  It is not a claim that $A$ and $T_A$ have the same truth value for
arbitrary independently chosen fresh-label values.  That distinction is
essential when using equisatisfiability in a proof checker.

\subsection{Resolution and Davis--Putnam elimination}

The resolution rule, with a variable $p$ absent from $D$ and $E$, is
\[
 \frac{D\cup\{p\}\qquad E\cup\{\neg p\}}{D\cup E}.
\]
The union removes duplicate literals.  A resolution certificate is a finite
list of input clauses and resolution steps, each step naming two earlier
indices and a pivot.  Its checker merely verifies those finite operations.

\begin{lemma}[Resolution soundness]\label{ca-resolution-sound}
Every model of the two parent clauses is a model of their resolvent.
Consequently a checked resolution certificate ending in the empty clause
certifies unsatisfiability.
\end{lemma}
\begin{proof}
If the resolvent were false, all literals in $D$ and $E$ would be false.
The positive parent would force $p=1$ while the negative parent would force
$p=0$, impossible for a bit.  Equivalently, split directly on the value of
$p$: one of the parent clauses must be satisfied by its remaining literals.
Induction along the certificate list preserves truth of every derived clause.
The empty clause has value $0$ at every assignment, proving the last claim.
\end{proof}

Partition a normalized clause set $F$ into clauses $F_0$ not mentioning $p$,
clauses $p\vee D_i$, and clauses $\neg p\vee E_j$.  Define
\[
 \operatorname{elim}_p(F)=F_0\cup\{D_i\cup E_j:i,j\},
\]
discarding tautologies and duplicate clauses after forming the resolvents.

\begin{lemma}[Elimination with a witness reconstruction rule]
\label{ca-elimination}
An assignment $b$ to the other variables satisfies
$\operatorname{elim}_p(F)$ if and only if it has an extension satisfying $F$.
The extension is computed as follows: if some $D_i$ is false at $b$, set
$p=1$; otherwise set $p=0$.
\end{lemma}
\begin{proof}
A model of $F$ satisfies $F_0$ and all resolvents by resolution soundness, so
its restriction is a model of the elimination.  Conversely suppose $b$
satisfies the elimination.  If every $D_i$ is true, setting $p=0$ satisfies
the positive clauses through their $D_i$ parts and satisfies all negative
clauses through $\neg p$.  If some $D_i$ is false, setting $p=1$ satisfies
all positive clauses.  For every $j$, the resolvent $D_i\cup E_j$ is true:
it is either retained in the elimination or was discarded as a tautology.
Since $D_i$ is false, $E_j$ is true.  Thus every negative clause is also true.
The clauses of $F_0$ hold in either case.  The selected case and, when needed,
the least offending $i$ are computed by a bounded scan.
\end{proof}

\begin{theorem}[Resolution completeness with primitive recursive bounds]
\label{ca-resolution-complete}
For every finite clause set $F$ on $n$ variables there is a primitive
recursive procedure returning either a satisfying assignment or a checked
resolution refutation.  After normalization the refutation can be represented
as a directed acyclic list with at most $(n+1)3^n$ lines, using no variables
outside the input vocabulary.
\end{theorem}
\begin{proof}
Eliminate variables in their fixed order, retaining each earlier clause and
recording a resolution justification for each new resolvent.  After $n$
rounds there are no variables.  The remaining normalized set is either empty,
which is satisfied by the empty assignment, or contains the empty clause.
These alternatives are decided by inspecting its finite code.
In the first case reconstruct the variable values in reverse order by
Lemma~\ref{ca-elimination}.  Each step gives a model of the preceding clause
set, ultimately a model of $F$.  In the second case, the empty clause already
has a finite list of input or resolution justifications; deleting discarded
clauses from the active set does not remove their earlier proof lines.
This gives a checked refutation.  There are at most $3^n$ normalized input
clauses, and each of the $n$ stages needs to store at most $3^n$ distinct new
clauses.  This proves the stated generous bound.  Duplicate candidates can
reuse an earlier stored line.

For totality, all clause sets are bitsets over the fixed universe of at most
$3^n$ clauses.  Each round scans at most $9^n$ parent pairs, each pivot and
each union is a bounded operation, and there are exactly $n$ rounds.  Proof
indices are bounded by the displayed line bound.  Hence the solver,
back-substitution, and proof recording are all primitive recursive.
Correctness is an induction on these fixed rounds, maintaining the bounded
equivalence in Lemma~\ref{ca-elimination} and the proof-checking invariant.
It does not presuppose an unbounded search for a refutation.
\end{proof}

\begin{corollary}[Clausal consequence by refutation]\label{ca-clausal-entail}
For a finite clause set $F$ and a clause
$C=\ell_1\vee\cdots\vee\ell_r$, the assertion $F\models_2 C$ is equivalent
to existence of a resolution refutation of
$F\cup\{\overline\ell_1,\ldots,\overline\ell_r\}$, where the added literals
are unit clauses.  A refutation can be constructed within the preceding
bound for the combined variable set.
\end{corollary}
\begin{proof}
A counterexample to $F\models_2C$ is exactly a model of the augmented set:
it satisfies $F$ and makes every literal of $C$ false.  Apply soundness and
completeness of the previous theorem.  This formulation deliberately uses
refutation; a pure resolution system without weakening need not directly
derive every syntactically larger consequence clause from $F$.
\end{proof}

\subsection{DPLL as a bounded, correct search}

For a partial assignment $\rho$, simplify $F$ by removing every clause already
made true and removing every false assigned literal from the remaining
clauses.  Denote the residual set by $F\!\upharpoonright\!\rho$.

\begin{lemma}[Restriction, unit propagation, and pure literals]
\label{ca-dpll-rules}
The following transformations preserve the indicated satisfiability problem.
\begin{enumerate}
 \item An extension of $\rho$ satisfies $F$ exactly when its unassigned part
       satisfies $F\!\upharpoonright\!\rho$.
 \item If the residual set contains a unit clause $\{\ell\}$, every model
       makes $\ell$ true; assigning it true loses no model.
 \item If a residual literal occurs and its complement does not, assigning
       that literal true preserves existence of a model.
 \item For an unassigned variable $p$, the residual set is satisfiable
       exactly when one of its restrictions $p=0,p=1$ is satisfiable.
\end{enumerate}
\end{lemma}
\begin{proof}
For the first claim, consider each clause separately.  A removed satisfied
clause imposes no further condition, and deleting its already false literals
leaves exactly the condition on its unassigned literals.  The second claim
is the definition of a unit disjunction.  For the third claim, take any
residual model and change the pure variable to make its occurring sign true.
Every clause mentioning it becomes or stays true; every other clause is
unchanged, since the opposite sign occurs nowhere.  Restricting the changed
assignment gives a residual model, and any residual model extends back by
the selected value using the first claim.  The final claim splits on the
two possible values of $p$ in a model; conversely a model of either restricted
set extends by that value.
\end{proof}

\begin{theorem}[Total and correct DPLL]\label{ca-dpll}
Starting with $n$ unassigned variables, the following deterministic policy
decides satisfiability by a primitive recursive search: report failure if an
empty clause occurs, report success if no clauses remain, otherwise assign
the least unit literal if one exists, then the least pure literal if one
exists, and otherwise split on the least occurring unassigned variable.
Every path assigns at most $n$ variables, and the full search tree contains
at most $2^{n+1}-1$ nodes.
\end{theorem}
\begin{proof}
Maintain the residual formula and the list of unassigned variables.
Whenever neither terminal case applies, some unassigned variable occurs;
each unit, pure, or split operation removes one such variable from that
list.  Define the search by primitive recursion on a fuel $r$ initially
equal to the list length.  Every recursive child receives fuel $r-1$.
At $r=0$ the only residual clauses are empty, so the terminal cases exhaust
the possibilities.  Every returned successful assignment is checked by
Lemma~\ref{ca-dpll-rules}, extending arbitrary unused variables by $0$ when
no clauses remain.  For completeness, the same lemma says that a model
survives a unit or pure step and must survive one child of a split.  Induction
on $r$ therefore proves both correctness assertions.  For the node bound,
$t(0)=1$ and $t(r+1)\leq1+2t(r)$ give $t(r)\leq2^{r+1}-1$;
unary steps only decrease this bound.  This recursive call tree can equally
be evaluated by a bounded explicit stack for that many steps, so no
general recursive termination principle is required.
\end{proof}

\subsection{Horn clauses: least models and forward certificates}

A Horn clause has at most one positive literal.  After deleting tautologies,
write it as a rule
\[
  p_{i_1}\wedge\cdots\wedge p_{i_k}\longrightarrow q,
\]
where $q$ is a variable or $\bot$.  Rules with $q=\bot$ are constraints;
rules with a variable head are definite.  The empty body represents $\top$.
For a set $Y$ of variables define
\[
 U(Y)=Y\cup\{q:\text{a definite input rule has head $q$ and body in $Y$}\}.
\]
Starting with $Y_0=\varnothing$, compute $Y_{t+1}=U(Y_t)$ for $n$ rounds.

\begin{theorem}[Horn least-model theorem and decision procedure]
\label{ca-horn}
The set $Y_n$ is closed under every definite input rule and is contained in
every model of those rules.  The full Horn set is satisfiable exactly when
no constraint has all its body variables in $Y_n$.  In that case $Y_n$ is
its least model under inclusion.  If a constraint does fire, forward chaining
provides a finite contradiction certificate.
\end{theorem}
\begin{proof}
Every $Y_t$ is contained in $Y_{t+1}$.  If equality occurs at one stage it
persists, since both next stages apply the same function $U$.  A strict
change adds at least one of the $n$ variables.  If $Y_n$ were not fixed,
there would have been $n+1$ strict increases starting from the empty set,
contradicting the count of distinct variables.  Thus $Y_n=U(Y_n)$, exactly
the required definite-rule closure.

Let $M$ be any model of the definite rules.  The empty set is contained in
$M$.  If $Y_t\subseteq M$ and a rule fires at $Y_t$, every body variable
is true in $M$, so its head is true in $M$.  Therefore
$Y_{t+1}\subseteq M$.  Induction gives $Y_n\subseteq M$.  If no constraint
fires at $Y_n$, definite closure and that fact verify every input clause,
so $Y_n$ itself is a model.  If a constraint fires, its body lies in every
possible model by the inclusion just proved, whereas its head is false;
thus no model exists.

For certificates record, when a variable is first added, one rule whose
body was already available at the preceding stage.  This produces an
acyclic list of at most $n$ variable conclusions, followed, in the failure
case, by one constraint conclusion $\bot$.  The checker verifies the input
rule and the earlier availability of its body variables.  If ordinary
resolution clauses are desired, resolve the corresponding Horn clause
successively with the unit clauses for those variables.  This derives its
head unit, or the empty clause for a fired constraint.  A body with $k$
literals uses $k$ resolution steps.  The resulting certificate length is
bounded by $(n+1)(L+1)$, where $L$ is the total input literal count.
\end{proof}

\begin{corollary}[Closure of Horn models under intersection]
\label{ca-horn-intersection}
The intersection of a nonempty finite list of models of a Horn set is again
a model.  For definite rules the same holds for the empty intersection,
interpreted as the set of all variables.
\end{corollary}
\begin{proof}
If a body is true in the intersection, it is true in each listed model.
A variable head is consequently true in each and hence in the intersection.
For a constraint, its body cannot be true in even one listed model, so it
cannot be true in the intersection.  With an empty list there are no
constraints in the definite case, and every head is true in the full set.
The qualification about a nonempty list cannot be dropped for a Horn set
containing, for example, the constraint $p\to\bot$.
\end{proof}

Forward chaining also has a worklist implementation.  Give each rule a
counter initially equal to its number of body occurrences, and for each
variable keep the list of occurrences in rule bodies.  Enqueue each
zero-counter head; when a variable first becomes true, process its occurrence
list and decrease the corresponding counters.  A counter reaching zero
enqueues its head or reports the fired constraint.  Mark true variables so
that duplicate enqueues are not processed again.  Each variable is processed
once and each body occurrence once, and each rule fires once.  Thus the
number of indexed array operations is $O(n+L+1)$, including initialization.
The invariant says that a counter equals the number of its not-yet-processed
body occurrences; it proves the same closure computation as above.
This operation count uses an indexed-array cost model.  The simpler fixed
$n$-round implementation already supplies everything required for an iPRA
formalization, without any machine-cost assumptions.

\subsection{Two-literal clauses and finite graph reachability}

Suppose every clause has at most two literals and no empty clause occurs.
Replace a unit clause $a$ by $a\vee a$.  Form a directed graph on the $2n$
literals; each clause $a\vee b$ gives edges
$\overline a\to b$ and $\overline b\to a$.  For every edge $u\to v$ the
complement-reversed edge $\overline v\to\overline u$ also occurs.

\begin{lemma}[Finite reachability and components]\label{ca-reachability}
Reachability on a graph with $m$ vertices is primitive recursive, and can be
computed by at most $m^3$ Boolean updates after initialization.  Mutual
reachability is an equivalence relation.  Its component graph has no directed
cycle between distinct components, and a topological ranking of its
components is computable by bounded search.
\end{lemma}
\begin{proof}
Initialize $R_0$ with edges and identity.  For $k=0,\ldots,m-1$, put
\[
 R_{k+1}(i,j)=R_k(i,j)\vee\bigl(R_k(i,k)\wedge R_k(k,j)\bigr).
\]
The invariant is that $R_k(i,j)$ holds exactly when there is a path from
$i$ to $j$ whose internal vertices have indices below $k$.  For a path using
$k$, remove loops until it is simple, and split it at its unique occurrence
of $k$; both resulting paths have internal vertices below $k$.  Conversely
concatenate the two paths witnessed by the second term.  Loop removal needs
at most the length of the supplied path, and a simple path has at most
$m-1$ edges.  Equivalently the invariant can carry these bounded path
witnesses explicitly.  This proves the final reachability table correct.
Reflexivity and transitivity follow by zero-length paths and concatenation,
and mutual reachability is therefore an equivalence relation.  A component
cycle would give paths both ways between its vertices, merging those
components, so it cannot occur.

To rank components, repeatedly remove the least remaining component with no
incoming edge from another remaining component.  Such a component exists:
otherwise choose the least predecessor at each step for one more step than
the number of remaining components.  Two positions in this finite list must
coincide, giving a cycle.  The pigeonhole claim is proved by marking distinct
components as they are encountered; after at most the component count all
marks have been used.  Thus bounded search always finds a source.  Assign
successive ranks in the removal order.  Every edge goes from a component of
no greater rank to one of no smaller rank.
\end{proof}

\begin{theorem}[The 2-SAT component criterion]\label{ca-twosat}
A two-literal clause set without an empty clause is satisfiable if and only
if no variable and its complement lie in the same strongly connected
component of the implication graph.  A satisfying assignment or a pair of
contradictory paths can be produced by a primitive recursive procedure.
\end{theorem}
\begin{proof}
Every model of a clause respects each associated implication edge: if
$\overline a$ is true in a model of $a\vee b$, then $a$ is false and $b$
is true.  Truth therefore propagates along paths.  If $p$ and $\neg p$ are
mutually reachable, split on the value of $p$.  Whichever of the pair is
true forces the other true along one of those paths, a contradiction.  The
reachability construction retains witnesses of length at most $2n-1$.

Conversely assume the component condition.  Complementation maps components
to components: reverse and complement each edge of a path, and mutual
reachability is preserved.  No component equals its complement.  Let $r(C)$
be a topological rank, and assign a literal $a$ to be true exactly when
\[
 r([a])>r([\overline a]).
\]
The ranks are distinct for complementary components, so exactly one of
$a,\overline a$ is true; the definition is a consistent variable assignment.
If both literals of a clause $a\vee b$ were false, its two edges would give
\[
 r([a])<r([\overline a])\leq r([b])
       <r([\overline b])\leq r([a]),
\]
which is impossible for natural numbers.  Hence every clause is satisfied.
The finite ranks and comparisons give the required explicit assignment.
The graph and its reachability table have size bounded in $n$ and the input
length, and all later searches are bounded by $2n$, establishing primitive
recursiveness.  An input empty clause is handled separately as immediate
unsatisfiability.
\end{proof}

\section{Forgetting variables, interpolation, and finite definability}
\label{ca-interpolation}

The quantifier notation introduced here names finite operations on formula
codes.  It does not add first-order or second-order quantifiers to the object
language.  Partition a formula's variables into a forgotten tuple
$P=(p_0,\ldots,p_{k-1})$ and a retained tuple $Q$.  Define
\[
 \exists_P A=\bigvee_{a<2^k} A[P:=a],
 \qquad
 \forall_P A=\bigwedge_{a<2^k} A[P:=a].
\]
The substituted values are $\bot,\top$.  Each formula contains only variables
of $Q$.  If $s$ is the syntax length of $A$, a size bound is
$(s+2)2^k$ with the usual binary folds and fixed constant symbols.

\begin{theorem}[Boolean elimination of finite quantifiers]
\label{ca-forgetting}
For every assignment $b$ to $Q$,
\[
 (\exists_P A)[b]=1\ \Longleftrightarrow\
        \text{some }a<2^k\text{ has }A[a,b]=1,
\]
\[
 (\forall_P A)[b]=1\ \Longleftrightarrow\
        \text{every }a<2^k\text{ has }A[a,b]=1.
\]
In the existential case the least witness is recovered by bounded search.
Both operations are primitive recursive.
\end{theorem}
\begin{proof}
Substitution commutes with evaluation, so the displayed formula evaluations
are exactly the disjunction and conjunction of the stated $2^k$ bits.
An induction on the fold length proves that a Boolean disjunction equals
$1$ exactly when some scanned bit equals $1$, recording the first such
position; the dual induction proves the conjunction assertion.  A fixed loop
over $2^k$ substitutions and a fixed fold construct the formula codes.
\end{proof}

\begin{theorem}[Uniform interpolation in the retained vocabulary]
\label{ca-uniform}
If $C$ contains no variable of $P$, then
\[
 A\models_2 C\quad\Longleftrightarrow\quad
           \exists_P A\models_2 C,
 \qquad
 C\models_2 A\quad\Longleftrightarrow\quad
           C\models_2\forall_P A.
\]
Thus $\exists_P A$ is the strongest Boolean consequence of $A$ expressible
without $P$, and $\forall_P A$ is the weakest Boolean sufficient condition
for $A$ expressible without $P$.
\end{theorem}
\begin{proof}
Certainly $A\models_2\exists_P A$: the current forgotten-variable assignment
selects one disjunct.  Thus a consequence of $\exists_P A$ is a consequence
of $A$.  Conversely suppose $A\models_2 C$ and a retained assignment makes
$\exists_P A$ true.  Extract a forgotten assignment making $A$ true and
apply the entailment.  Since $C$ contains no forgotten variables, its value
does not change when that witness is forgotten.  This proves the first
equivalence, even if $C$ contains extra variables disjoint from $P$; those
extra values are held fixed throughout.

For the second equivalence, $\forall_P A\models_2 A$ is immediate by
selecting the current conjunct.  In the other direction, if $C$ is true and
$C\models_2 A$, vary the forgotten assignment through every $a<2^k$.
The value of $C$ stays true, so every cofactor of $A$ is true; hence the
conjunction $\forall_P A$ is true.  The characterizations ``strongest'' and
``weakest'' are just these two universal entailment properties.
\end{proof}

\begin{corollary}[Finite propositional interpolation]\label{ca-craig}
If $A\models_2 B$, there is an explicitly computed formula $I$ using only
variables common to $A$ and $B$ such that
$A\models_2 I$ and $I\models_2 B$.  One may take
\[
 I=\exists_{\operatorname{Var}(A)\setminus\operatorname{Var}(B)} A.
\]
By finite classical completeness, corresponding checked classical proofs of
the two implications can also be constructed.
\end{corollary}
\begin{proof}
The forgotten variables do not occur in $B$, so
Theorem~\ref{ca-uniform} gives $I\models_2 B$ and the first part of its
proof gives $A\models_2 I$.  Every variable surviving in $I$ occurs in both
formulas.  All variable-set operations are finite scans of the input syntax.
Apply the truth-table proof constructor of Theorem~\ref{thm:cpc-complete} to each
of the two verified implications.  This last step is an additional proof
construction and is not needed for semantic interpolation itself.
If the common vocabulary is empty, $I$ is a constant formula, as required:
it is true precisely when some assignment satisfies $A$.
\end{proof}

\begin{proposition}[Laws of forgetting]\label{ca-forgetting-laws}
For disjoint forgotten tuples $P,R$, finite quantifier operations commute
and compose:
\[
 \exists_P\exists_R A\equiv_2\exists_{P\cup R}A
 \equiv_2\exists_R\exists_P A,
\]
with the analogous assertion for $\forall$.  Forgetting twice is idempotent.
Moreover
\[
 \exists_P(A\vee B)\equiv_2(\exists_PA)\vee(\exists_PB),\qquad
 \forall_P(A\wedge B)\equiv_2(\forall_PA)\wedge(\forall_PB),
\]
\[
 \forall_P A\equiv_2\neg\exists_P\neg A.
\]
If $P$ is absent from $C$, then
\[
 \exists_P(A\wedge C)\equiv_2(\exists_PA)\wedge C,
 \qquad
 \forall_P(A\vee C)\equiv_2(\forall_PA)\vee C.
\]
\end{proposition}
\begin{proof}
The composition statements identify an assignment to the union with a pair
of assignments to the separate tuples, simply a reindexing of a finite
Boolean fold.  Idempotence holds because after the first forgetting the
cofactors for the forgotten variables are identical.  The distribution
identities express associativity and commutativity of finite joins or meets.
Duality follows by complementing each bit in a finite conjunction.  For the
last two identities, $C$ has a fixed bit across all cofactors; split on that
bit.  For example, when $C=0$ both sides of the existential identity vanish,
and when $C=1$ both are $\exists_PA$.  These arguments also cover an empty
forgotten tuple, which has exactly one assignment.  Existential forgetting
does not in general preserve conjunction: forgetting $p$ from $p\wedge\neg p$
gives $\bot$, whereas $(\exists_p p)\wedge(\exists_p\neg p)$ is $\top$.
\end{proof}

\begin{theorem}[Finite Beth definability]\label{ca-beth}
Let $T(Q,P)$ be a finite propositional theory, written as one conjunction,
and let $p$ be a member of $P$.  Suppose the following finite uniqueness
condition holds: any two assignments satisfying $T$ and agreeing on $Q$
agree on $p$.  Then the formula
\[
 d(Q)=\exists_P\bigl(T(Q,P)\wedge p\bigr)
\]
explicitly defines $p$ in the theory:
\[
 T\models_2\bigl(p\leftrightarrow d(Q)\bigr).
\]
The uniqueness condition is a bounded truth-table test, and the definition
and a classical proof certificate can be computed primitively recursively
whenever that test succeeds.
\end{theorem}
\begin{proof}
Take a model $(b,a)$ of $T$.  If its $p$ bit is $1$, that very assignment
is a witness to the finite existential defining $d$, so $d(b)=1$.  If
$d(b)=1$, extract a witness $a'$ making $T(b,a')\wedge p$ true.  The two
models $(b,a)$ and $(b,a')$ agree on $Q$, so the supplied uniqueness condition
forces their $p$ bits to agree; the second is $1$, hence so is the first.
This proves both directions of the biconditional at every model.  A retained
assignment with no extension to a model of $T$ simply makes $d$ false; no
existence-of-model assumption was used.  Check the uniqueness condition by
scanning all triples $(b,a,a')$, a finite bounded table.  Theorem~\ref{ca-forgetting}
constructs $d$ and extracts its witnesses.  Finally use the classical
truth-table proof constructor of Theorem~\ref{thm:cpc-complete} on the
established implication.  Repeating
the construction for every member of $P$ gives an explicitly defined tuple
whenever the whole tuple is uniquely determined by $Q$.
\end{proof}

\subsection{Arithmetic audit of the constructions}
\label{ca-audit}

The constructions in these sections use truth tables of length $2^n$, clause
universes of size $3^n$, and, for exact DNF minimization, subsets of a set of
size $3^n$.  Their bounds may be large, but exponentiation and finite
iteration of these explicit operations are primitive recursive.  A program
does not stop merely because a mathematical argument says that an unbounded
search eventually succeeds: it runs a fixed number of rounds, retaining its
answer once a stopping condition has been met.  The relevant bounds are
$n$ for elimination, Horn saturation, and variable assignment along a DPLL
path; $2^{n+1}-1$ for the complete DPLL tree; $m$ for the stages of finite
reachability; and the displayed table bounds for representation and
forgetting.  Numerical encodings of the finite lists may enlarge these bounds
by further primitive recursive functions without changing the conclusion.

For formal proofs in iPRA, an assertion about all assignments is stored as
the conjunction of its bounded Boolean checks.  A loop invariant says that
this Boolean invariant equals $1$.  Induction on the loop counter is therefore
quantifier-free induction in the language with the requisite primitive
recursive functions.  Witness-producing assertions retain the actual
assignment, path, formula, or certificate code in the loop state; correctness
of that retained code is again a bounded check.  A uniformly named Boolean
table supplied as an input code is distinct from an arbitrary higher-order
function.  These observations are the specific reason the preceding
arguments fit iPRA, beyond the bare fact that their algorithms terminate.

The breadth of the results should not be confused with a sharp complexity
classification.  The chapter proves the displayed bounds, and the stated
worklist operation count, but does not use or establish $\mathrm P\ne
\mathrm{NP}$, lower bounds for general propositional proof systems, or
compactness for arbitrary infinite premise sets.  Every interpolation or
definability claim in this chapter is explicitly finite and Boolean.
